\scanpage{206}
\mspage{205}

\hfill Tribu n° 59

DÉCISIONS sur les ALGÈBRES de LIE SEMI-SIMPLES

(n° 400)

-:-:-:-

Le nouveau titre sera ``algèbres de Lie semi-simples déployées'' ; les autres seront dites ``repliées''. Le prochain rédacteur suivra d'assez près la présente rédaction et y ajoutera :

a) Les modules infinis $W(\lambda)$, (cf. § 4), avec formules de Harish-Chandra, et résultats dudit sur les opérateurs différentiels invariants ;

b) Jacobson-Mort aux Off ;

c) Les algèbres déployées classiques, et leurs représentations fondamentales, (ne parler que des représentations d'algèbres) ; montrer les autom. de trialité de $D_4$ ; les algèbres exceptionnelles iront en exercices.

D'autre part Cartier rédigera un tapis arithmétique sur les $Z$-structures. Les autres compléments dont parlent les commentaires du n° 400 sont, ou bien vidés, ou bien incorporés dans les §§ déjà rédigés.

\uline{Annexe} I \uline{des applications polynomiales}.

Noter que l'annexe II est traitée avec la rédaction des racines.

\uline{n° 1} -- Dans l'énoncé de la prop. 1 ajouter ``En particulier toutou vert non vide est dense'', (ce qui supprime le cor.) ; prendre la démonstration par références.

\uline{n° 2} -- A la p. A3, l.7, ajouter ``pour les topol. de Zariski''. Juste après intervertir définition et critère pour les applic. dominantes. A la l.7 bas (p. A3), remplacer $W$ par $A_W$.
